\section{Overview of \barrel}

This section presents \barrel and the workflow depicted in \cref{fig:barrel-archi}.
We first recall the B method, then explain how \barrel handles well-definedness conditions, and finally detail its inner workings and automation features.

\subsection{Background}

\paragraph*{The B method and proof obligations.}
The B method is a state-based formal development approach centered on \emph{abstract machines} and stepwise \emph{refinement}.
A machine declares abstract sets and constants, a state described by variables constrained by an inductive invariant, an initialization, and operations.
A development proceeds by refining an abstract machine into more concrete refinements down to an implementation, introducing representation data together with a \emph{gluing invariant} relating concrete and abstract states.
From each machine and refinement step, Atelier~B generates \emph{proof obligations} (POs): logical sequents expressing the required invariant-preservation and refinement conditions.
Concretely, POs cover invariant preservation by the initialization and each operation under its precondition, and simulation-style obligations ensuring that concrete operations preserve the abstract behavior through the gluing invariant.
In addition, the PO generator emits \emph{well-definedness} (WD) obligations for the many partial constructs of the language.

POs can be discharged automatically using Atelier~B’s internal provers and by delegating to external backends.
In particular, several approaches translate B obligations to SMT solvers such as Z3~\cite{z3} or cvc5~\cite{cvc5} via dedicated encodings, including the \emph{ppTransSMT}~\cite{DeharbeBSMT} and the more recent \emph{BEer}~\cite{beer} encodings.
Obligations that remain---often those requiring nontrivial quantifier instantiation, refinement reasoning, or additional WD reasoning---are typically discharged interactively using Atelier~B’s GUI-oriented prover, where the user guides rule applications and instantiations.

\paragraph*{The B language in a nutshell.}
% Set theory lies at the core of the B method.
B's mathematical language is based on classical set theory and first-order logic with equality.
Specifications are written using the usual set-theoretic constructors such as comprehension, powerset, Cartesian product, and an algebra of relations.
Ordered pairs are first-class values (written as \emph{maplets} $x \mapsto y$), and relations are simply sets of such pairs, i.e.,\ $R \subseteq A \times B$.
Functions are not primitive: a partial function $f$ from $A$ to $B$ (written $f \in A \pfun B$) is a relation $f \subseteq A \times B$ whose domain is contained in $A$ and which is \emph{functional}:
for all $x \in A$ and $y,z \in B$, if $(x \mapsto y) \in f$ and $(x \mapsto z) \in f$ then $y = z$.
Total ($A \longrightarrow B$), injective (partial $A \pinj B$ and total $A \tinj B$), and surjective (partial $A \psurj B$ and total $A \tsurj B$) variants are defined similarly.

% That is, every value is a set that is constrained by some properties that define its shape.
% For example, a (total) function $f$ from a set $A$ to a set $B$ (written $f \in A \longrightarrow B$) is a set of pairs in $A \times B$ satisfying the functionality property: for every $x \in A$ and $y,z \in B$, if $(x, y) \in f$ and $(x, z) \in f$ then $y = z$.
% However, specifications written in the B method are still required to satisfy some typing predicates \todo[perhaps let's not talk about ``typing'' B?]: a ``parent set'' must be deducible for every constant and variable of a specification.
Although the underlying logic is set-theoretic and thus untyped, Atelier~B enforces a static well-formedness discipline often described as ``typing'': every identifier and expression must have a deducible ``parent set.''
This guarantees in particular that constructed sets are homogeneous and rules out ill-formed terms.
These deduced parent sets provide exactly the information needed to interpret B expressions in Lean's typed logic in our translation.

\begin{example}
	The expression $1 + \top$ is well-formed in ZFC set theory---and definitionally equal, under the encoding \(\bot \hookrightarrow 0\) and \(\top \hookrightarrow 1\), to the set $\mathbb{B}$ of Booleans, or the natural number $2$---but it is not a valid expression in B.
	Indeed, $1$ has parent set $\mathbb{Z}$ and $\top$ has parent set $\mathbb{B}$, and $+$ expects both its arguments to have parent set $\mathbb{Z}$.

\end{example}

\paragraph*{Foundations of B in Lean.}
%
We use these typing annotations to ground our encoding in typed higher-order logic (rather than ZFC set theory), using Mathlib's \texttt{Set} and \texttt{SetRel} types\footnote{These are defined, respectively, as $\mathtt{Set}\ \alpha \triangleq \alpha \to \mathtt{Prop}$ and $\mathtt{SetRel}\ \alpha\ \beta \triangleq \mathtt{Set}(\alpha \times \beta)$.} to represent sets and relations:
each inferred B carrier is mapped to a corresponding Lean type---\texttt{INTEGER} to \texttt{Int}, Cartesian products to Lean product types, and powersets to \texttt{Set} over the element type.
This is a convenient representation since Mathlib comes
% batteries included\fixme[remove?] 
with many definitions and theorems related to \texttt{Set} for us to use when proving the translated obligations in Lean.
This representation also allows us to express every kind of set used in B specifications since all sets in B are required to be \emph{homogeneous}, i.e.,\ restricted to a uniform shape of elements.
For instance, the set $\{0, \text{``hi''}\}$ must be considered ill-formed by implementations of the B method.
% this also enables us to easily reason about sets using the plethora of theorems and definitions that Mathlib puts at our disposal.

We start by defining standard B operators directly in Lean---reusing their notations from B, in order for the POs to be easily understandable by B users---and typing them appropriately in Lean to keep their use as generic as possible.
Some examples (the set of relations, the set of partial functions, the set of total injective functions, and domain subtraction) are given in \cref{lst:B-sets-lean}.
%
\begin{lstlisting}[label=lst:B-sets-lean, caption={Encoding of some B operators in Lean.}, float=t]
abbrev rels {α β} (A : Set α) (B : Set β) : Set (SetRel α$$ β) :=
  𝒫$$ (A ×$$ˢ B)
scoped infixl:125 "  ⟷$$  " =>$$ rels

abbrev pfun {α β} (A : Set α) (B : Set β) : Set (SetRel α$$ β) :=
  { f ∈$$ A  ⟷$$  B | ∀$$ ⦃x y z⦄, (x, y) ∈$$ f →$$ (x, z) ∈$$ f →$$ y = z }
scoped infixl:125 " ⇸$$ " =>$$ pfun

abbrev injTFun {α β} (A : Set α) (B : Set β) : Set (SetRel α$$ β) :=
  A ⤔$$ B ∩$$ A  ⟶$$  B
scoped infixl:125 " ↣$$ " =>$$ injTFun

abbrev domSubtr {α β} (E : Set α) (R : SetRel α$$ β) : SetRel α$$ β$$ :=
  { z ∈$$ R | z.1 ∉$$ E }
scoped infixl:160 " ⩤$$ " =>$$ domSubtr
\end{lstlisting}
%
We then provide proofs, using Mathlib's infrastructure and theorems, of various laws and theorems of B: for any total function $f \in A \longrightarrow B$ its domain $\bdom{f}$ equals $A$; assuming that $a \le b$, the minimum of $\binterval{a}{b}$ is $a$; etc.
Both the definitions and theorems are readily available to end users when importing \barrel in a Lean project, and serve as basic building blocks for proofs of B obligations translated by our tool.



\subsection{Encoding partial B operators}
\label{sec:partial-operators}
A central challenge in this setting is the use of \emph{partial} B operators, meaning that they are only defined when some \emph{well-definedness} (WD) conditions are satisfied.
A typical example is the minimum operator $\bmin{S}$, which returns the smallest element of a set $S$ (of integers or reals) with respect to $\le$, if it exists, and is undefined otherwise.
These partial operators are typically encoded using Hilbert's bounded choice operator $\epsilon\ x \in S.\ P$, which is left unspecified if no element of $S$ satisfies $P$.
The \emph{WD condition} of the corresponding operator is precisely that $P$ is satisfied by at least one element of $S$.



Atelier~B generates separate WD obligations from a component's source clauses.
The principal POs do not carry these obligations as explicit proof arguments; establishing the development's correctness therefore requires relating the separate WD checks to the principal proofs.
This generation does not cover arbitrary terms introduced subsequently by interactive commands.
In the example in \cref{sec:lean-ip-comparison}, an interactive instantiation leads the prover to accept a contradictory assertion.

In \barrel, dependencies between WD goals and principal goals are represented by explicit proof arguments, as in the encoding of Russell~\cite{sozeau:2006} into Rocq.
%
First, we generate WD conditions to match the exact context in which they are inserted.
Atelier~B's generated WD conditions may differ because they undergo simplification and because our encoding uses different definitions for some predicates and WD conditions.



\begin{remark}
	\barrel defines the WD of minimum as the existence of a least element in the set, and maximum dually as the existence of a greatest element.
	The definitions apply uniformly to partially ordered types.
	For integers, the library proves equivalence with the non-emptiness and finiteness conditions used by Atelier~B: a minimum requires a finite negative part, and a maximum a finite non-negative part.
	For reals, mere non-emptiness and boundedness would not suffice: the interval $(0,1)$ has a lower bound but no least element.
	The integer equivalences are stated as follows:
	\begin{lstlisting}[language=lean]
theorem min.WD_iff_AtelierB_WD {S : Set ℤ} :
	min.WD S ↔$$ S ≠$$ ∅$$ ∧$$ S ∩$$ (INTEGER \ NATURAL) ∈$$ FIN INTEGER
theorem max.WD_iff_AtelierB_WD {S : Set ℤ} :
	max.WD S ↔$$ S ≠$$ ∅$$ ∧$$ S ∩$$ NATURAL ∈$$ FIN INTEGER
\end{lstlisting}
\end{remark}

These equivalences relate the operator conditions; they do not imply a one-to-one correspondence between the generated WD obligations.
Atelier~B may simplify its WD goals and remove hypotheses, whereas \barrel generates conditions in the exact Lean context of each encoded occurrence and shares their proofs when possible.
% \footnote{\todo[Should this footnote be kept?] To be accurate, our WD conditions are generated from the exact Lean context whereas Atelier~B already performs simplification steps at this point, which may alter the PO.}
%
Then, we use Lean's dependent types to encode partial operators using \texttt{Classical.choose}\footnote{Hilbert's $\epsilon$ operator, whose type is \mbox{$\{\alpha : \mathtt{Sort}\ u\} \to \{p : \alpha \to \mathtt{Prop}\} \to (\exists\, x,\, p\ x) \to \alpha$}}, which explicitly depends on a proof that the predicate is satisfied by at least one element.\footnote{The encoding represents partial operators as total Lean functions with explicit WD proof parameters.}
This means that one cannot construct a term depending on a partial operator without providing a proof that it is well-defined.
The statement of this proof, which is the WD condition of the operator, is then generated as a side theorem (or a subgoal) which must be discharged for the main theorem to be fully proven, i.e.,\ not depend on the \texttt{sorryAx} axiom---the axiom that is introduced by admitted theorems.
To illustrate this use of \texttt{Classical.choose}, the encoding (and WD conditions) of $\bmin{S}$ and function application $\bapp{F}{x}$ are given in \cref{lst:min-fn-app}.

\begin{lstlisting}[language=lean, mathescape=true, label=lst:min-fn-app, caption={Encoding of the $\bmin{S}$ and $\bapp{F}{x}$ partial operators.}, float=t]
structure min.WD {α} [PartialOrder α] (S : Set α) : Prop where
  isBoundedBelow : ∃$$ x ∈$$ S, ∀$$ y ∈$$ S, x ≤$$ y

noncomputable 
abbrev min {α} [PartialOrder α] (S : Set α) (wd : min.WD S) : α$$ :=
  Classical.choose wd.isBoundedBelow

structure app.WD {α β} (f : SetRel α$$ β) (x : α) : Prop where
  isPartialFunction : f ∈$$ dom f  ⟶$$  ran f
  isInDomain : x ∈$$ dom f -- $\textcolor{commentcolor}{\equiv\ \exists}$ y, (x, y) ∈$$ f

noncomputable 
abbrev app {α β} (f : SetRel α$$ β) (x : α) (wd : app.WD f x) : β$$ :=
  Classical.choose wd.isInDomain
\end{lstlisting}
%$$

\subsection{Inner workings of \barrel and automation}
\label{sec:innner-workings}

\paragraph*{User-level workflow.}

\barrel is organized into four distinct passes.
First, the machine, system, refinement, or implementation is parsed into the BXML format, and proof obligations are generated from this format using Atelier~B's internal tools \texttt{bxml} and \texttt{pog}.

Although Atelier~B's \texttt{pog} can generate WD conditions for the given machine, we choose to rely on our own generation instead, since the WD conditions that are generated by \texttt{pog} may not exactly match the context in which they will appear in the Lean terms.
Then, the resulting POs (in XML format) are parsed by \barrel into an abstract syntax tree (AST) and normalized.
Normalization ensures that B predicates like $P \land Q \Rightarrow R$ are treated the same as the logically equivalent predicate $P \Rightarrow Q \Rightarrow R$.
Normalization is not required for the translation, but reduces the number of conjunction-elimination steps needed in the resulting Lean goals.
The normalized B formulas are then transcribed into Lean terms using Lean's facilities for metaprogramming and elaboration.
Encoding partial B operators requires tracking the WD conditions inserted during translation.
We use Lean's metavariables, which retain their local contexts, to represent these conditions.
A WD metavariable is assigned an application of a theorem with the required type: either a suitable theorem already available for reuse (via exact type matching or subsumption), or a fresh obligation extracted from its local context.
Subsumption, here, allows deriving a fact from a \texttt{@[subsume]}-tagged theorem with fewer hypotheses (i.e.,\ a more general statement), e.g.,\ deriving \(P \Rightarrow Q \Rightarrow R\) from \(P \Rightarrow R\), \(Q \Rightarrow R\), or just \(R\).
Proved WD theorems can be reused within an import, but also, once the component is completed, by subsequent imports, including across Lean modules; \cref{sec:proof-workflow} explains the smaller theorem interfaces that make this reuse possible across refinements.
%
\begin{example}
	%
	We illustrate how \barrel handles WD conditions using the following (already normalized) B predicate, stating that every nonempty finite set of integers---i.e.,\ belonging to $\bfinone{\mathbb{Z}}$---whose elements are nonpositive has a nonpositive minimum: \[
		\forall S \cdot (S \in \bfinone{\mathbb{Z}} \Rightarrow (\forall x \cdot x \in S \Rightarrow x \le 0) \Rightarrow \bmin{S} \le 0)
	\]
	%
	The translation first generates a Lean term while inserting metavariables in place of WD conditions (where $h_1$, $h_2$, and $h_3$ are fresh identifiers):
	\[
		\begin{split}
			 & P :=                                                                                                                                                                    \\
			 & \quad \forall\, (S : \mathtt{Set}\, \mathbb{Z})\, (h_1 : S \in \mathtt{B.Builtins.FIN_1}\ \mathbb{Z})\, (h_2 : \forall\, (x : \mathbb{Z})\, (h_3 : x \in S),\ x \le 0), \\
			 & \qquad\qquad \mathtt{B.Builtins.min}\ S\ (\metavar{m_1}\ S\ h_1\ h_2) \le 0
		\end{split}
	\]
	The new metavariable $\metavar{m_1}$ is generated internally by \barrel and contains $S : \mathtt{Set}\ \mathbb{Z}$, $h_1 : S \in \mathtt{B.Builtins.FIN_1}\ \mathbb{Z}$, and $h_2 : \forall\, (x : \mathbb{Z})\ (h_3 : x \in S),\ x \le 0$ in its local context.
	Although its type is $\mathtt{B.Builtins.min.WD}\ S$, its local context is implicitly universally quantified, hence the full application $\metavar{m_1}\ S\ h_1\ h_2$ in the goal.
	%
	A fresh theorem statement $min\_wd_1$ is then generated---and remains to be proved---from the type of $\metavar{m_1}$ universally quantified over its local context: \[
		\begin{split}
			 & \mathtt{theorem}\ min\_wd_1 :                                                                                                                                          \\
			 & \quad \forall\, (S : \mathtt{Set}\ \mathbb{Z})\, (h_1 : S \in \mathtt{B.Builtins.FIN_1}\ \mathbb{Z})\, (h_2 : \forall\, (x : \mathbb{Z})\, (h_3 : x \in S),\ x \le 0), \\
			 & \qquad\quad \mathtt{B.Builtins.min.WD}\ S
		\end{split}
	\]
	$\metavar{m_1}$ is finally assigned $min\_wd_1$, and substituted in $P$ to obtain the final goal to be proved by the user: \[
		\begin{split}
			\vdash \forall\, & (S : \mathtt{Set}\ \mathbb{Z})\, (h_1 : S \in \mathtt{B.Builtins.FIN_1}\ \mathbb{Z})\, (h_2 : \forall\, (x : \mathbb{Z})\, (h_3 : x \in S),\ x \le 0), \\
			                 & \quad \mathtt{B.Builtins.min}\ S\ (min\_wd_1\ S\ h_1\ h_2) \le 0
		\end{split}
	\]
\end{example}

These steps are performed by the \texttt{import <type> <name> from <folder>} command of \barrel.
The input format is specified by \texttt{<type>}: \texttt{machine}, \texttt{system}, \texttt{refinement}, \texttt{implementation}, or \texttt{pog}.
The argument \texttt{<name>} gives the file name without the extension.
For POG files supplied directly, the conversion to BXML and PO generation are skipped.
Each remaining goal can then be selected with \texttt{next obligation} or, by name, with \texttt{obligation <po>}.
Both commands accept \texttt{from <term>} or \texttt{by <proof>} to prove WD conditions and reconstructed proof obligations.
Finally, the import must be closed with \texttt{qed <name>}, which checks that all obligations have been discharged, and only then adds them to the global context as theorems.
However, \texttt{qed} does not check for dependencies on \texttt{sorryAx}; users must check these themselves, as in other Lean developments.

% : \begin{itemize}
% 	\item 
% 	\item 
% 	\item Then, the formulas are transcribed into Lean terms using Lean's facilities for meta-programming/elaboration.
% 	\item Finally, these Lean terms and their associated WD conditions are presented to the user who fills out the proofs.
% \end{itemize}

% The first three steps are performed by the command \texttt{import machine <name> from <folder>}.

% \noindent\todo[Inner workings]

\paragraph*{Proof automation.}
\label{par:automation}
Before presenting any goals to the user, \barrel runs an automation pass whose primary purpose is to discharge routine side conditions induced by our encoding of B partial operators.
Concretely, the tactic layer repeatedly applies relevant lemmas---backtracking as needed---that relate WD subgoals and available hypotheses.
These helper lemmas provide definitional equations and rewriting and simplification facts for the supported B constructs involving sets, relations, functions, arithmetic, and finiteness. Many generated WD goals then reduce to elementary facts about membership, domain/range inclusion, finiteness, and order bounds and can be closed automatically.
We use Mathlib's facilities to tag helper lemmas with relevant categories---\texttt{wd\_min}, \texttt{wd\_max}, \texttt{wd\_app}, \texttt{wd\_card}, \texttt{wd\_mod}, \texttt{wd\_div}---and its set-theoretic results to discharge the corresponding goals.
These tags also name the tactics run automatically by \barrel's automation layer; users may invoke them manually or extend them with their own lemmas.
For function application, the WD tactic can use a total-function membership lemma tagged \texttt{@[tfun]} together with domain membership; \cref{ex:workflow-update} illustrates this extension on a function update in the case study.




Beyond the core encoding of B primitives, \barrel provides a library of theorems about built-in operators.
Its purpose is twofold:
\begin{enumerate}
	\item provide reusable WD facts that discharge the side goals generated for partial operators;
	\item offer canonical rewrite lemmas that turn B-shaped expressions into ``Mathlib-friendly'' goals.
\end{enumerate}
For instance, in the arithmetic fragment used throughout our running examples, the library includes dedicated WD lemmas for the minimum operation, as shown in \cref{lst:min-wd-lemmas}, each of them being tagged with the relevant theorem category.

\begin{lstlisting}[language=lean, caption={WD lemmas about extrema in \barrel.}, label=lst:min-wd-lemmas, float]
@[wd_min]
theorem min.WD_singleton {α : Type _} [PartialOrder α] {a : α} :
	min.WD {a} :=$$ ⋯
@[wd_min]
theorem min.WD_of_finite {α : Type _} [LinearOrder α] {S A : Set α}
	(h : S ∈ FIN₁ A) : min.WD S :=$$ ⋯
@[wd_min]
theorem min.WD_interval {lo hi : ℤ} (h : lo ≤$$ hi) :
	min.WD (lo..hi) :=$$ ⋯
\end{lstlisting}

The library also provides membership results about finite sets and intervals, together with computational equalities tagged as rewrite rules for automation (\cref{lst:finset-interval-lemmas}). Proofs are replaced by $\leanpf$ for brevity.

\begin{lstlisting}[language=lean, caption={Membership lemmas about finite sets and intervals in \barrel.}, label=lst:finset-interval-lemmas, float]
theorem interval.FIN₁_mem {lo hi : ℤ} (h : lo ≤$$ hi) :
	lo .. hi ∈$$ FIN₁ INTEGER :=$$ ⋯
@[simp]
theorem min.of_singleton {α : Type _} [PartialOrder α] {a : α} :
	min {a} (min.WD_singleton) = a :=$$ ⋯
@[simp]
theorem interval.min_eq {lo hi : Int} (h : lo ≤$$ hi) :
	min (lo .. hi) (min.WD_interval h) = lo :=$$ ⋯
\end{lstlisting}

The remaining obligations---typically invariant preservation and refinement simulation goals---are left as interactive proof tasks.